\documentclass[notes]{beamer}
% \documentclass{beamer}

\usepackage[utf8]{inputenc}
\usepackage[T1]{fontenc}

\usepackage{pgfpages}

\usepackage{listings}
\usepackage{listings-rust}
\usepackage{color}
\usepackage{verbatim}

\usepackage{siunitx}

\usepackage[frenchb]{babel}
\usepackage[left=\flqq, right=\frqq, leftsub=\flqq, rightsub=\frqq]{dirtytalk}

\usetheme{Warsaw}

% \setbeameroption{show notes on second screen}

\makeatletter
\newenvironment{withoutheadline}{
  \setbeamertemplate{headline}[default]
  \def\beamer@entrycode{\vspace*{-\headheight}}
}{}
\makeatother

\setbeamertemplate{footline}
{%
  \leavevmode
  \begin{beamercolorbox}[wd=.5\paperwidth,ht=2.5ex,dp=1.125ex,leftskip=.3cm plus1fill,rightskip=.3cm]{section in head/foot}%
    Vincent BONNEVALLE % TODO mettez vos nom
  \end{beamercolorbox}%
  \begin{beamercolorbox}[wd=.5\paperwidth,ht=2.5ex,dp=1.125ex,leftskip=.3cm,rightskip=.3cm plus1fil]{subsection in head/foot}%
    Présentation du langage Rust \hfill \Tiny{\textbf{\color{white}{\insertframenumber/\inserttotalframenumber}}}
  \end{beamercolorbox}%
}

\begin{document}

\title{Un titre intelligent sur Rust}
\author{Julien, Hugo, Vincent Bonnevalle} % Mettez vos nom :)
\institute{Université Paris-Didierot}
\date{}
\titlegraphic{\includegraphics[width=.25\textwidth,height=.25\textheight]{logo.pdf}}

\begin{withoutheadline}
  \begin{frame}
    \titlepage{}
  \end{frame}
\end{withoutheadline}

  \section{Introduction}
  \begin{frame}
    \begin{itemize}
    \item zero-cost abstraction
    \item move semantics
    \item guaranted memory safety
    \item program without data races
    \item trait-based generics
    \item pattern matching
    \item type inference
    \item minimal runtime
    \item efficient C bindings
    \end{itemize}
  \end{frame}

\begin{withoutheadline}
  \begin{frame}
    \frametitle{Plan}
    \setcounter{tocdepth}{1}
    \tableofcontents{}
    \setcounter{tocdepth}{2}
  \end{frame}
\end{withoutheadline}

\section{Section}
\begin{frame}
  \tableofcontents[currentsection, hideothersubsections]
\end{frame}

\subsection{Du code}
\begin{frame}[fragile]
  Contexte d'unification :
\begin{lstlisting}[language=Rust]
let c = "Du code !!!"
\end{lstlisting}
\end{frame}
\subsection{Séparation de la diapo en 2}
\begin{frame}
  \begin{columns}
    \begin{column}{0.5\textwidth}
      Une colonne
    \end{column}
    \begin{column}{0.5\textwidth}
      l'autre colonne
    \end{column}
  \end{columns}
\end{frame}

\section{Syntaxe de Rust}
\begin{frame}
\end{frame}
\section{Le typer de Rust}
\begin{frame}
\end{frame}
\section{Rust et la parallélisation}
\begin{frame}
\end{frame}

\end{document}
